% ============================================================================
% 30_ai.tex  --  AI-collaboration note  (~1/3 page)
% Two threads: (i) generation + correction of the physics, (ii) using the model
% to judge what is essential at a fixed page budget, and to learn the register
% of scientific writing.
% ============================================================================

\section{On the AI collaboration}
\label{sec:ai}

An LLM mined the block notes and reproduced the chiral-basis determinant
argument of Sec.~\ref{sec:majoranalines} from Ref.~\cite{leumer2021thesis} in
the convention of Eq.~\eqref{eq:HK}. Three corrections came out of that
exchange: the measured ratios (Table~\ref{tab:fits}) forced the
$(1-\lambda^2)$ prefactor onto an initially bare $2t|\lambda|^L$; a physical
reading of the large-$L$ spikes was dropped once Eq.~\eqref{eq:splitting} was
seen to predict the opposite trend; and an adversarial re-check caught a sign
slip in Eq.~\eqref{eq:char} that had propagated into $\lambda$.

Compressing the material to the page limit proved the more instructive half. What
to cut is a physics judgement as much as an editorial one --- a step earns its
place only if the argument cannot be reconstructed without it --- and settling
that forced an explicit view of what each section is for. Presentation is a
technique in its own right, to be practised like the physics: redrafting
against the model sharpened the register, showing where a hedge is honest and
where it is padding, and how a scientific sentence is built to carry one claim
at a time. Every number was re-verified against the code; the draft is to be
revised by the author.
